\documentclass[aspectratio=169]{beamer}
\usetheme{metropolis}
\usepackage{nexus_theme}
\usepackage{listings}

\title{Designing Your Research Question}
\subtitle{Module 2.1: Finding Literature}
\author{Nexus Scholar Suite --- Modern Research Methodologies}
\date{}

\begin{document}

% --- TITLE SLIDE ---
\begin{frame}[plain]
  \titlepage
\end{frame}

% --- LEARNING OBJECTIVES ---
\begin{frame}{What You'll Learn}
  By the end of this lesson, you will be able to:
  \begin{enumerate}
    \item \textbf{Construct} a research question using PICO, SPIDER, or Heilmeier.
    \item \textbf{Distinguish} between a topic, a gap, and a question.
    \item \textbf{Perform} a systematic gap analysis.
    \item \textbf{Validate} your question using the FINER criteria.
  \end{enumerate}
\end{frame}

% --- THE FUNNEL ---
\begin{frame}{Topic $\neq$ Gap $\neq$ Question}
  \begin{center}
  \resizebox{0.8\linewidth}{!}{
  \begin{tikzpicture}
    \node[card, fill=gray!10, draw=gray, minimum width=12cm, minimum height=1.2cm] (topic) {\textbf{TOPIC:} ``Machine learning in healthcare''};
    \node[primarycard, minimum width=9cm, minimum height=1.2cm, below=0.3cm of topic] (gap) {\textbf{GAP:} ``No study has compared transformers to RNNs for pediatric ICU readmission''};
    \node[successcard, minimum width=6cm, minimum height=1.2cm, below=0.3cm of gap] (rq) {\textbf{QUESTION:} ``How does transformer-based prediction compare to RNN-based models in AUROC?''};
    
    \draw[->, thick, gray] (topic.south) -- (gap.north);
    \draw[->, thick, gray] (gap.south) -- (rq.north);
  \end{tikzpicture}
  }
  \end{center}
  \vspace{0.2cm}
  \textit{Most PhD students confuse a topic for a question. This leads to years of unfocused work.}
\end{frame}

% --- FRAMEWORKS ---
\begin{frame}{Frameworks for Crafting Questions}
  \begin{columns}[T]
    \begin{column}{0.32\textwidth}
      \begin{block}{\textbf{PICO}\\(Health Sciences)}
        \begin{itemize}
          \item \textbf{P}opulation
          \item \textbf{I}ntervention
          \item \textbf{C}omparison
          \item \textbf{O}utcome
        \end{itemize}
      \end{block}
    \end{column}
    \begin{column}{0.32\textwidth}
      \begin{exampleblock}{\textbf{SPIDER}\\(Qualitative)}
        \begin{itemize}
          \item \textbf{S}ample
          \item \textbf{P}henomenon
          \item \textbf{D}esign
          \item \textbf{E}valuation
          \item \textbf{R}esearch type
        \end{itemize}
      \end{exampleblock}
    \end{column}
    \begin{column}{0.32\textwidth}
      \begin{alertblock}{\textbf{Heilmeier}\\(Engineering/CS)}
        \begin{itemize}
          \item What are you doing?
          \item How is it done today?
          \item What's new?
          \item Who cares?
          \item What are the risks?
        \end{itemize}
      \end{alertblock}
    \end{column}
  \end{columns}
\end{frame}

% --- GAP ANALYSIS ---
\begin{frame}{Systematic Gap Analysis}
  \begin{enumerate}
    \item \textbf{Search broadly:} Use \texttt{scholar-search-kit} with a wide query.
    \item \textbf{Build the citation graph:} Use \texttt{scholar-graph-kit} to visualize connections.
    \item \textbf{Find white space:} Look for disconnected or weakly-connected clusters in the network. These are under-explored intersections.
    \item \textbf{Read ``Future Work'':} The highest-PageRank papers' Future Work sections are literal roadmaps to publishable questions.
  \end{enumerate}
\end{frame}

% --- FINER CRITERIA ---
\begin{frame}{The FINER Test}
  Before committing to your question, validate it:
  \vspace{0.3cm}
  \begin{itemize}
    \item[\ding{51}] \textbf{F}easible --- Can I answer this with my resources and time?
    \item[\ding{51}] \textbf{I}nteresting --- Will my committee and the field care?
    \item[\ding{51}] \textbf{N}ovel --- Has this exact question been answered?
    \item[\ding{51}] \textbf{E}thical --- Can I pursue it without causing harm?
    \item[\ding{51}] \textbf{R}elevant --- Does the answer advance knowledge?
  \end{itemize}
  \vspace{0.3cm}
  \textit{If it fails even one criterion, refine before investing months of work.}
\end{frame}

% --- EXERCISE PREVIEW ---
\begin{frame}{Your Turn: Exercise}
  \begin{exampleblock}{Exercise: The Question Forge}
    Starting from a broad topic, narrow it down through a gap analysis and produce a FINER-validated research question using the appropriate framework for your discipline.
  \end{exampleblock}
\end{frame}

% --- STANDOUT CLOSE ---
\begin{frame}[standout]
  A good research question is the difference\\
  between a thesis that takes three years\\
  and one that takes seven.
\end{frame}

\end{document}
